% Pressure-scope revision: analyses/021-2026-09-10-training-results-figures/pressure-scope/observations/O016-pressure-placement.md; evidence.py, figures.py, tables.py.
% Revision/source crosswalk: analyses/018-2026-09-08-results-materials/observations/O014-manuscript-polish.md
% Evidence: Analysis 018 observations O001/O002/O003/O004/O007/O011.
% OL1 budget regimes: Analysis 021 observations/O002-ol1-geometry.md; 02_ol1_geometry.py.
% Two-size extension: analyses/024-2026-09-17-h-only-kernel-latency/observations/025-matched-quality-latency.md;
% generating analysis: analyses/024-2026-09-17-h-only-kernel-latency/16_plot_matched_quality_latency.py.
% Cross-size scope and A0 histories: Analysis 021 observations/O006-scale-transfer.md
% and O007-a0-optimization.md; 05_scale_transfer.py and 06_a0_optimization.py.
\section{Discussion and conclusion}
\label{sec:conclusion}

% Scale extension: analyses/038-2026-09-24-base-t2-t7-scale-frontier/observations/003-manuscript-integration.md; 04_manuscript_evidence.py.
Our results distinguish the benefits of targeted thresholding from the costs of broader pressure. At 14M, execution ablations and optimization diagnostics identify $h,z$ as useful sparsification sites under the studied implementation. Comparisons of $T_2/P_h$ and $T_7/P_h$ across 14M, 31M and 70M further show that narrowing threshold coverage preserves validation quality at matched thresholds, while broader thresholding can reach lower latency at a larger quality cost. The validation-loss improvement over the dense control is confined to the evaluated 14M conditions. The broader result is a favorable quality--latency trade-off from targeting useful sites, rather than a general improvement in quality from sparsification.

More broadly, sparsification is one instance of a training--inference co-design problem. The relevant question is not simply where a model can be made sparse, but which representation makes expensive computation removable at inference time. A systems-guided approach is therefore to identify the operations that dominate latency, determine the structured property that would make them cheaper or skippable, and enforce that property during training. This principle may extend beyond activation sparsity to tiled activation or weight sparsity, quantization, batched execution, and other hardware-dependent forms of structure.

\subsection{Failures and Limitations}

Our main analysis combines extensive 14M intervention ablations with targeted-versus-broad thresholding comparisons at 14M, 31M and 70M. We also trained 410M models, but the fixed-token budget places that size in a different optimization regime: with the same token budget, the 410M base model has higher loss than 70M and larger gradient norms (Appendix~\ref{app:410m-results}). Kernel efficiency also does not transfer directly across scales. Adapting the 14M kernel to 70M introduces substantial overhead; further optimization improves execution, but does not close the gap with native PyTorch, although it benefits from the sparse structure. Our runtime study is deliberately narrow (batch-one, full-sequence inference on a single GPU), aimed to provide a controlled setting for studying the interaction between train-time sparsification and execution, rather than a complete production-serving evaluation. Training comparisons use one seed per condition. This limits claims about individual point estimates, although the consistent trade-off structure across dozens of matched recipes provides evidence for the broader trends we report. Finally, the intervention space remains only partially explored. The 31M comparison holds pressure placement fixed at $h$; it does not independently establish the spillover mechanism, site-wise optimization conflict or conditional execution savings measured at 14M. Pressure and threshold placements are combinatorial, and broad-pressure experiments change several pressure targets at the same time. Because OL1 also averages pressure across targeted sites, these comparisons do not isolate the contribution of each added site. Future work could separate pressure targets one at a time, remove this cross-site averaging confound, and allow site-specific thresholds rather than sharing a single $\kappa$ across sites.

